% Coxeter s. 88 

\begin{proposition}[\pol{twierdzenie o dwusiecznej}\ita{teorema della bisettrice}]
	\pol{Dany jest trójkąt $\triangle ABC$.}
	\ita{Sia dato un triangolo $\triangle ABC$.}
	\pol{Odcinek $CF$, gdzie $F$ leży na odcinku $AB$, jest dwusieczną kąta przy wierzchołku $C$ wtedy i tylko wtedy, gdy zachodzi równość:}
	\ita{Il segmento $CF$, dove $F$ appartiene al segmento $AB$, è la bisettrice dell'angolo al vertice $C$ se e solo se vale l'uguaglianza}
	\begin{equation}
		\frac{|AF|}{|BF|} = \frac{|AC|}{|BC|}.
	\end{equation}
	\pol{\index{twierdzenie!o dwusiecznej}}
	\ita{\index{teorema!della bisettrice}}
\end{proposition}

\pol{To jest (VI.3) w Elementach Euklidesa.}
\ita{Questa è la proposizione (VI.3) degli \emph{Elementi} di Euclide.}
\pol{Dowód korzysta z podobieństwa trójkątów, twierdzenia sinusów albo własności pola, patrz}
\ita{La dimostrazione utilizza la similitudine dei triangoli, il teorema dei seni oppure proprietà delle aree; si veda}
\pol{\index{twierdzenie!sinusów}}%
\ita{\index{teorema!dei seni}}%
Guzicki \pol{\cite[s. 120]{guzicki_2021}} \ita{\cite[p. 120]{guzicki_2021}} \pol{(siedem różnych dowodów);}
\ita{(sette dimostrazioni differenti);}
Bogdańska, Neugebauer \pol{\cite[s. 73]{neugebauer_2018},}
\ita{\cite[p. 73]{neugebauer_2018},}
Audin \pol{\cite[s. 102]{audin_2003}} \ita{\cite[p. 102]{audin_2003}} \pol{w formie ćwiczenia;}
\ita{come esercizio;}
\pol{patrz też \cite[s. 11]{komisarczyk_2024}.}
\ita{si veda anche \cite[p. 11]{komisarczyk_2024}.}
Hartshorne \pol{\cite[s. 181]{hartshorne2000}} \ita{\cite[p. 181]{hartshorne2000}} \pol{zostawia je jako ćwiczenie dla Czytelnika.}
\ita{lo lascia come esercizio al lettore.}

\pol{Uogólnienie do kątów dwuściennych w czworościanie znajdzie Kieza w $\Delta_{12}^3$.}
\ita{Una generalizzazione agli angoli diedri di un tetraedro è presentata da Kieza in $\Delta_{12}^3$.}
% https://www.deltami.edu.pl/2012/03/odkryj-twierdzenie-sam/

\begin{proposition} % Guzicki s. 126
	\pol{Punkt $P$ znajdujący się wewnątrz kąta wypukłego leży na dwusiecznej tego kąta wtedy i tylko wtedy, gdy odległości tego punktu od ramion kąta są równe.}
	\ita{Un punto $P$ situato all'interno di un angolo convesso appartiene alla sua bisettrice se e solo se le sue distanze dai lati dell'angolo sono uguali.}
\end{proposition}

%

Wnioskiem z twierdzenia o dwusiecznej jest: % lang-pl
Una conseguenza del teorema della bisettrice è la seguente: % lang-it

\begin{theorem}
[Steinera-Lehmusa] % lang-pl
[di Steiner-Lehmus] % lang-it
\index{twierdzenie!Steinera-Lehmusa}% % lang-pl
\index{teorema!di Steiner-Lehmus}% % lang-it
\label{theorem_steiner_lehmus}%
    Jeżeli dwie dwusieczne trójkąta są równej długości, to trójkąt ten jest równoramienny. % lang-pl
	Se due bisettrici di un triangolo hanno la stessa lunghezza, allora il triangolo è isoscele. % lang-it
\end{theorem}

Po raz pierwszy wspomni o nim Christian Lehmus w~liście z~1840 roku do Charlesa Sturma, gdzie poprosi o~czysto geometryczny dowód. % lang-pl
Il teorema fu menzionato per la prima volta da Christian Lehmus in una lettera del 1840 a Charles Sturm, nella quale chiedeva una dimostrazione puramente geometrica. % lang-it
\index[persons]{Lehmus, Christian}%
\index[persons]{Sturm, Charles}%
Sturm przekazał prośbę do innych matematyków, jedną z~pierwszych osób, która uporała się z problemem, był Jakob Steiner. % lang-pl
Sturm trasmise la richiesta ad altri matematici; una delle prime persone a risolvere il problema fu Jakob Steiner. % lang-it
\index[persons]{Steiner, Jakob}%
Większość znanych uzasadnień to dowody nie wprost: jeśli trójkąt nie jest równoramienny, to ma dwusieczne różnej długości. % lang-pl
La maggior parte delle dimostrazioni note procede per assurdo: se un triangolo non è isoscele, allora le sue bisettrici hanno lunghezze diverse. % lang-it

Dowód można znaleźć u Hartshorne'a \cite[s. 11]{hartshorne2000}; Pompego \cite[s. 32, 33]{pompe_2022}; Bogdańskiej, Neugebauera \cite[s. 74]{neugebauer_2018}. % lang-pl
Zdzisław Pogoda w $\Delta_{11}^1$ napisze, jakie pomysłowe rozwiązanie znaleźli dwaj angielscy inżynierowie, G. Gilbert oraz D. McDonnell. % lang-pl
Una dimostrazione si trova in Hartshorne \cite[p. 11]{hartshorne2000}, Pompe \cite[pp. 32, 33]{pompe_2022} e Bogdańska, Neugebauer \cite[p. 74]{neugebauer_2018}. % lang-it
Zdzisław Pogoda, in $\Delta_{11}^1$, descrive l'ingegnosa soluzione trovata da due ingegneri inglesi, G. Gilbert e D. McDonnell. % lang-it
\index[persons]{McDonnell, D.}%
\index[persons]{Gilbert, G.}%
% TODO The American Mathematical Monthly (7,1963,str.79–80)
Coxeter \cite[s. 32]{coxeter_1967} podaje je w formie ćwiczenia (po tym jak wcześniej \cite[s. 26, 33]{coxeter_1967} poprosi o~dowód tego samego dla dwusiecznej zamienionej na środkową lub wysokość, co znacząco obniża poziom trudności). % lang-pl
Eves \cite[s. 19, 58]{eves1_1972} zachęca do dowodu z dopiskiem, że jest trudny. % lang-pl
Coxeter \cite[p. 32]{coxeter_1967} la propone come esercizio (dopo aver chiesto in precedenza \cite[pp. 26, 33]{coxeter_1967} una dimostrazione analoga con la bisettrice sostituita da una mediana o da un'altezza, il che riduce notevolmente la difficoltà). % lang-it
Eves \cite[pp. 19, 58]{eves1_1972} invita il lettore a dimostrarlo, osservando che non è facile. % lang-it

% https://www.algebra.com/algebra/homework/word/geometry/Medians-in-an-isosceles-triangle.lesson

%

% TODO: https://zadania.info/d801/3107027

\pol{Twierdzenie o symetralnej}
\ita{Il teorema dell'asse del segmento}
% TODO czy to jest oficjalna nazwa
\pol{\index{twierdzenie!o symetralnej}}
\ita{\index{teorema!dell'asse del segmento}}
\pol{można wysłowić tak: miejsce geometryczne punktów $X$, dla których $|AX|/|BX| = 1$, jest prostą.}
\ita{può essere formulato così: il luogo geometrico dei punti $X$ tali che $|AX|/|BX| = 1$ è una retta.}
\pol{Około dwusetnego roku przed naszą erą Apoloniusz z Pergi udowodni piękne uogólnienie tego faktu:}
\ita{Intorno al 200 a.C. Apollonio di Perga dimostrò una bella generalizzazione di questo fatto:}

\begin{definition}[\pol{okrąg Apoloniusza}\ita{circonferenza di Apollonio}] % TODO: Guzicki s. 129
\label{definicja_okregu_apoloniusza}%
	\pol{Dane są dwa różne punkty $A$, $B$ oraz liczba dodatnia $\lambda \neq 1$.}
	\ita{Siano dati due punti distinti $A$, $B$ e un numero positivo $\lambda \neq 1$.}
	\pol{Wtedy zbiór punktów}
	\ita{Allora l'insieme dei punti}
	\begin{equation}
		A(\lambda) := \left\{X : \frac{|AX|}{|BX|} = \lambda \right\}
	\end{equation}
	\pol{jest okręgiem o środku na prostej $AB$ i promieniu równym}
	\ita{è una circonferenza con centro sulla retta $AB$ e raggio pari a}
	\begin{equation}
		R = \frac{\lambda}{|\lambda^2 - 1|} \cdot |AB|.
	\end{equation}
	\pol{\index{okrąg!Apoloniusza}}
	\ita{\index{circonferenza!di Apollonio}}
\end{definition}

% TODO: https://www.deltami.edu.pl/2013/01/okrag-apoloniusza/

Coxeter \pol{\cite[s. 104, 105]{coxeter_1967}} \ita{\cite[pp. 104, 105]{coxeter_1967}} \pol{pisze, że ten okrąg odbija inwersyjnie punkty $A, B$ (a dla $\lambda = 1$ dostajemy symetralną odcinka).}
\ita{scrive che questa circonferenza scambia inversivamente i punti $A$ e $B$ (e per $\lambda = 1$ si ottiene l'asse del segmento).}
\pol{Wspomni o nim Eves \cite[s. 159]{eves1_1972}}
\ita{Ne parla anche Eves \cite[p. 159]{eves1_1972}.}

\pol{To jest jeden z pięciu okręgów Apoloniusza, oprócz tego mamy dwie rodziny wzajemnie prostopadłych okręgów, okręgi Apoloniusza trójkąta (pomocne w znajdowaniu punktów izodynamicznych oraz prostej Lemoine'a), okrąg z problemu Apoloniusza styczny do trzech danych oraz fraktal zwany po angielsku uszczelką -- \emph{,,Apollonian gasket''}.}
\ita{Questa è una delle cinque nozioni chiamate circonferenze di Apollonio; vi sono inoltre due famiglie di circonferenze mutuamente ortogonali, le circonferenze di Apollonio di un triangolo (utili per trovare i punti isodinamici e la retta di Lemoine), la circonferenza del problema di Apollonio tangente a tre oggetti dati e il frattale noto in inglese come \emph{``Apollonian gasket''}.}
% https://en.wikipedia.org/wiki/Circles_of_Apollonius
\todofoot{\pol{Loksodromiczny ciąg Coxetera okręgów}\ita{Sequenza lossodromica di circonferenze di Coxeter}} % https://en.wikipedia.org/wiki/Coxeter%27s_loxodromic_sequence_of_tangent_circles

\pol{Piszą o nim Bogdańska, Neugebauer \cite[s. 74]{neugebauer_2018}}
\ita{Ne scrivono Bogdańska e Neugebauer \cite[p. 74]{neugebauer_2018}.}

%